% Generated draft: requires mathematical and source review.
\documentclass[11pt]{article}
\usepackage[margin=25mm]{geometry}
\usepackage{fontspec}
\setmainfont{Latin Modern Roman}
\usepackage{amsmath,amssymb,mathtools,unicode-math}
\setmathfont{Latin Modern Math}
\usepackage{enumitem,needspace,xurl,hyperref}
\setlength{\parindent}{0pt}
\setlength{\parskip}{5pt}
\setlength{\emergencystretch}{3em}
\newcommand{\N}{\mathbb{N}}
\newcommand{\Z}{\mathbb{Z}}
\newcommand{\Q}{\mathbb{Q}}
\newcommand{\R}{\mathbb{R}}
\newcommand{\C}{\mathbb{C}}
\newcommand{\F}{\mathbb{F}}
\newcommand{\A}{\mathbb{A}}
\newcommand{\PP}{\mathbb{P}}
\newcommand{\E}{\mathbb{E}}
\newcommand{\eps}{\varepsilon}
\newcommand{\dd}{\,\mathrm{d}}
\begin{document}
\section*{UnsolvedMath ID: 701768}
\textbf{Source problem code:} catalogue entry
\subsection{Definitions and mathematical statement}
A group \(G\) is residually finite if, for every \(g\in G\setminus\{1\}\), there is a homomorphism \(\varphi:G\to F\) to a finite group such that \(\varphi(g)\ne1\). An infinite group is just infinite if every nontrivial normal subgroup has finite index. It is hereditarily just infinite if every subgroup of finite index is just infinite. These definitions do not impose finite generation; whether the source intends finite generation as part of its terminology is unresolved.

A group is amenable if it admits a left-translation-invariant finitely additive probability measure on all its subsets. It is solvable if its derived series reaches the trivial subgroup after finitely many steps. The elementary amenable groups form the smallest class containing all finite groups and all abelian groups and closed under subgroups, quotients, extensions, and directed unions.

The supplied problem asks four separate existence questions: does there exist a residually finite hereditarily just infinite group \(G\) satisfying each respective condition below? The group may differ between parts.
\begin{enumerate}
\item[(a)] \(G\) is amenable and not solvable.
\item[(b)] \(G\) is amenable and not elementary amenable.
\item[(c)] \(G\) has intermediate word growth.
\item[(d)] \(G\) has intermediate subgroup growth.
\end{enumerate}

For part (c), the usual word-growth interpretation requires \(G\) to be finitely generated. For a finite symmetric generating set \(S\), let
\[
\gamma_{G,S}(n)=\bigl|\{g\in G:|g|_S\le n\}\bigr|,
\qquad n\in\Z_{\ge1},
\]
where \(|g|_S\) is the least length of a word in \(S\) representing \(g\). Intermediate word growth means subexponential growth without a polynomial upper bound:
\[
\lim_{n\to\infty}\frac{\log\gamma_{G,S}(n)}{n}=0,
\qquad
\forall C>0\ \forall d\ge0\ \exists n\ge1:
\gamma_{G,S}(n)>Cn^d.
\]
This condition is independent of the finite generating set. Finite generation in this interpretation is made explicit here because it is not stated separately in the supplied question.

For part (d), define
\[
s_G(n)=\bigl|\{H\le G:[G:H]\le n\}\bigr|.
\]
The source does not define intermediate subgroup growth or explicitly require finite generation. One conventional interpretation, for finitely generated \(G\), is
\[
\lim_{n\to\infty}\frac{\log s_G(n)}{n}=0,
\qquad
\forall C>0\ \forall d\ge0\ \exists n\ge1:
 s_G(n)>Cn^d.
\]
Thus, under that interpretation, part (d) asks for a finitely generated example with subgroup growth faster than every polynomial upper bound and subexponential in the index bound. This interpretation is conditional: the supplied record does not determine whether this convention and finite-generation hypothesis are intended. Without finite generation or another finiteness hypothesis, the counts \(s_G(n)\) need not be finite.

The supplied contextual comment states: ``All examples that we know are either linear (hence Tits Alternative applies), or have a quotient with property (T) (hence cannot be amenable).'' This is preserved as a source-relative report about known examples, not as an independently verified classification. Its nonamenability inference requires an infinite quotient with property \((T)\); infinitude is not stated in the record. Finite quotients with property \((T)\) do not give that obstruction. The comment also leaves the hypotheses for its intended use of the Tits alternative implicit.
\subsection{Short English statement}
M. Vannacci asks whether residually finite hereditarily just infinite groups can be amenable but nonsolvable, amenable but not elementary amenable, of intermediate word growth, or of intermediate subgroup growth. These are four separate existence questions. Finite-generation conventions and the precise meaning of intermediate subgroup growth require clarification. The supplied registry labels the problem open; that label has not been independently verified.
\subsection{Sources}
\begin{itemize}
\item Local UnsolvedMath registry, supplied source record with permanent ID \(701768\), exact source code \texttt{catalogue entry}, title \emph{Kourovka Notebook Problem 20.110}. Supplied status: open. The record was created and updated on \texttt{2026-10-16T14:00:00Z}; these are registry timestamps, not literature-verification dates. Supplied upstream revision: \texttt{d371ee917ac9647fbb06cd6fc529deec6f3dda82}.
\item The supplied attribution is to M. Vannacci, proposed in 2022, in \emph{The Kourovka Notebook (Unsolved Problems in Group Theory)}, 20th issue (2022), original problem number 20.110. The background also identifies the 21st issue edition (2026), \texttt{arXiv:1401.0300v48}. The registry describes the collection as Issues 1--21 (1965--2026), edited by E. I. Khukhro and V. D. Mazurov. These bibliographic details and the contextual comment are supplied source information and have not been independently checked.
\end{itemize}
\end{document}
